% from ega3-6-fr.tex
\oldpage[III]{173}
\textnormal{(iii)} Dans la démonstration de (6.10.1), on peut supposer les
$M_{\bullet\bullet\bullet}^{(i)}$ formés de $A_i$-modules \emph{libres} (comme il
résulte aisément de la démonstration de (\textbf{0}, 11.5.2.1) « dualisée »); les
$M_{\bullet\bullet\bullet}^{(i)}\otimes_{A_i}B$
sont alors formés de $B$-modules libres, et comme
$M_{\bullet\bullet\bullet}$ est égal à leur produit tensoriel sur $B$, on voit
qu'on peut supposer en outre dans (6.10.1) que $\mathscr R_\bullet$ est associé
à un complexe de $B$-modules \emph{libres}. En outre, en vertu de (M, XVII, 1.2), le
tricomplexe $M_{\bullet\bullet\bullet}$ dépend \emph{fonctoriellement} de chacun des
bicomplexes $L_{\bullet\bullet}^{(i)}$ (donc de chacun des
$\mathscr P_\bullet^{(i)}$, lorsqu'on fixe un recouvrement fini de chacun des
$X^{(i)}$), les « morphismes » de tricomplexes devant être ici entendus comme
les \emph{classes d'homomorphismes} pour la relation d'homotopie; d'ailleurs, le
remplacement d'un recouvrement de $X^{(i)}$ par un recouvrement plus fin
donnant lieu pour les $L_{\bullet\bullet}^{(i)}$ à des homomorphismes définis
précisément à homotopie près (6.6.8), on voit finalement qu'avec la convention
précédente pour les morphismes, le tricomplexe
$M_{\bullet\bullet\bullet}$ est un \emph{foncteur} en chacun des
$\mathscr P_\bullet^{(i)}$. Nous préciserons cette dépendance fonctorielle, et
notamment le comportement de $\mathscr R_\bullet$ relativement à des suites
exactes $0\longrightarrow\mathscr P_\bullet^{\prime(i)}\longrightarrow\mathscr P_\bullet^{(i)}\longrightarrow\mathscr P_\bullet^{\prime\prime(i)}\longrightarrow0$
de complexes, dans le chapitre consacré à une algèbre générale de foncteurs
cohomologiques, mentionné dans (6.1.3).
\end{remarks}

\begin{scholium}[6.10.4]
Le fait que $\mathscr R_\bullet$ est formé de $\mathscr O_Y$-Modules
$S$-\emph{plats} entraîne aisément que $\mathscr H_\bullet(\mathscr R_\bullet\otimes_S\mathscr Q'_\bullet)$
est un foncteur \emph{homologique} en $\mathscr Q'_\bullet$ (voir le raisonnement de
(7.7.1)). C'est cette propriété qui, ainsi qu'on l'a mentionné en (6.1.1), est
la motivation de l'introduction de l'hypertor. Posons en effet
\[
  X=X_1\times_SX_2\times\ldots\times_SX_n,
  \qquad
  f=f_1\times_Sf_2\times_S\ldots\times_Sf_n,
\]
\[
  \mathscr P_\bullet
  =
  \mathscr P_\bullet^{(1)}
  \otimes_S\mathscr P_\bullet^{(2)}
  \otimes\ldots\otimes_S\mathscr P_\bullet^{(n)},
\]
\[
  X'=X\times_SS',
  \qquad
  Y'=Y\times_SS',
  \qquad
  f'=f\times_S1_{S'};
\]
les problèmes de changement de base amènent à étudier l'hypercohomologie $\mathscr H^\bullet_{f'}(\mathscr P_\bullet\otimes_S\mathscr N')$
en tant que foncteur par rapport au $\mathscr O_{S'}$-Module quasi-cohérent
$\mathscr N'$, ou encore l'hypercohomologie $\mathscr H^\bullet_f(\mathscr P_\bullet\otimes_S\mathscr N)$
comme foncteur en le $\mathscr O_S$-Module quasi-cohérent $\mathscr N$. Lorsque
les $\mathscr P_\bullet^{(i)}$ (donc aussi $\mathscr P_\bullet$) sont
$S$-\emph{plats}, il résulte de ce qui précède et de (6.7.6) que ce foncteur est bien
un foncteur \emph{cohomologique} en $\mathscr N$; mais il n'en est plus de même
lorsqu'on ne fait plus l'hypothèse de platitude sur les
$\mathscr P_\bullet^{(i)}$, et l'on ne peut plus alors aborder l'étude de $\mathscr H^\bullet_f(\mathscr P_\bullet\otimes_S\mathscr N)$
par les méthodes usuelles de l'Algèbre homologique.

Nous aurons toutefois surtout à utiliser le cas où $n=1$, $Y=S$ et où
$\mathscr P_\bullet$ est formé de $\mathscr O_X$-Modules $Y$-\emph{plats}. On a dans
ce cas le
\end{scholium}

\begin{theorem}[6.10.5]
Soient $Y=\operatorname{Spec}(A)$ un schéma affine noethérien,
$f:X\to Y$ un morphisme propre, $\mathscr P_\bullet$ un complexe de
$\mathscr O_X$-Modules cohérents, \textup{plats} sur $Y$, limité inférieurement. Il
existe alors un complexe $\mathscr L_\bullet$ de $\mathscr O_Y$-Modules,
limité inférieurement, dont les termes $\mathscr L_i$ sont des
$\mathscr O_Y$-Modules de la forme $\mathscr O_Y^{mi}$, et un isomorphisme
\[
  \mathscr H^\bullet
  (f,\mathscr P_\bullet\otimes_Y\mathscr Q_\bullet)
  \xrightarrow{\sim}
  \mathscr H_\bullet
  (\mathscr L_\bullet\otimes_Y\mathscr Q_\bullet)
  \tag{6.10.5.1}\label{III.6.10.5.1.fr}
\]
de $\partial$-foncteurs en le complexe $\mathscr Q_\bullet$ de
$\mathscr O_Y$-Modules quasi-cohérents, limité inférieurement. En outre, pour
tout morphisme $u:Y'\to Y$, on a, en posant
\[
  X'=X_{(Y')},
  \qquad
  f'=f_{(Y')},
  \qquad
  \mathscr P'_\bullet
  =
  \mathscr P_\bullet\otimes_Y\mathscr O_{Y'},
  \qquad
  \mathscr L'_\bullet=u^*(\mathscr L_\bullet)
\]
